\chapter{Threshold-Robust Surrogate Gradient (TrSG)}
\label{ch:trsg}

This chapter presents the second contribution of the thesis: \textbf{TrSG}, a surrogate-gradient (SG) formulation whose gradient magnitude is invariant to the threshold scale and whose active region adapts naturally with $V_{\text{thr}}$. \Cref{sec:trsg_background} reviews surrogate-gradient descent. \Cref{sec:trsg_two_families} dissects two existing families—Absolute-Scale (AS-SG) and Relative-Scale (RS-SG)—and shows how each becomes unstable when $V_{\text{thr}}$ is trainable. \Cref{sec:trsg_method} derives TrSG, and \Cref{sec:trsg_grad_analysis,sec:trsg_compatibility} provide a detailed gradient analysis with respect to $V_{\text{thr}}$ and $\tau$ together with compatibility checks across surrogate shapes.

\begin{figure}[t!]
    \centering
    \includegraphics[width=0.99\linewidth]{fig/second_main.pdf}
    \caption{\textbf{AS-SG vs.\ RS-SG vs.\ TrSG with a rectangular surrogate gradient ($\gamma=1$).}
    The curve depicts the membrane-potential distribution; colored boxes represent the surrogate-gradient region (horizontal extent: active gradient window; vertical extent: gradient magnitude).
    Panels (a) and (b) compare small ($V_{\text{thr}}^{l}\!\ll\!1$) and large ($V_{\text{thr}}^{l}\!\gg\!1$) thresholds.
    \emph{AS-SG} uses a fixed window of width $\gamma$, leading to \textit{gradient flooding} or \textit{starvation}.
    \emph{RS-SG} normalizes the window by $V_{\text{thr}}^{l}$ but scales the magnitude as $1/V_{\text{thr}}^{l}$, causing \textit{explosion} or \textit{vanishing}.
    \textbf{TrSG} multiplies the threshold on the forward path during training, canceling the $1/V_{\text{thr}}^{l}$ factor and keeping both window and magnitude balanced.}
    \label{fig:second_main}
\end{figure}

\section{Motivation}
\label{sec:trsg_motivation}

Surrogate-gradient (SG) descent~\cite{wu:2018} enables gradient-based optimization of SNNs by replacing the non-differentiable spiking function with a differentiable surrogate near the threshold. This permits effective gradient descent and often improves performance at early timesteps. However, our analysis shows that SG used for internal neuronal parameters~\cite{fang:2021,wang:2022,rathi:2023}—particularly $V_{\text{thr}}$—makes training highly sensitive to the \emph{scale} of $V_{\text{thr}}$. While careful tuning can yield strong accuracy, this dependency exposes a fundamental limitation of current SG-based training.

In this chapter, we identify—\emph{to our knowledge for the first time}—that prior SG designs exhibit strong sensitivity of gradient flow to $V_{\text{thr}}$. To address it, we introduce \textbf{TrSG (Threshold-Robust SG)}, which maintains stable gradients regardless of the threshold value. TrSG is simple to implement, adds negligible overhead, and is compatible with a wide range of surrogate functions.

\section{Background: Surrogate Gradient Descent}
\label{sec:trsg_background}

By \cref{eq:lif_discrete_firing}, an LIF neuron emits a spike when the membrane potential crosses the threshold. Let $H(x)$ denote the Heaviside step. The spike can be written as
\[
S^{l}[t] \;=\; H\!\bigl(x(M^{l}[t],V^{l}_{\text{thr}})\bigr),
\]
where the \emph{argument} $x(\cdot)$ encodes how the membrane potential $M^{l}[t]$ and the threshold $V^{l}_{\text{thr}}$ are combined (e.g., $x = M^{l}[t]-V^{l}_{\text{thr}}$). Because $H(\cdot)$ is non-differentiable and $H'(\cdot)$ is a Dirac delta, SG training replaces $H$ with a smooth surrogate $f$ whose derivative locally approximates $H'$:
\[
H'(x)\;\approx\; f'(x)\quad\text{near }x=0.
\]
Applying the chain rule yields the working gradient
\begin{equation}
\label{eq:def_of_srgd}
\frac{\partial S^{l}[t]}{\partial M^{l}[t]}
\;=\;
H'\!\bigl(x\bigr)\,\frac{\partial x}{\partial M^{l}[t]}
\;\approx\;
f'\!\bigl(x\bigr)\,\frac{\partial x}{\partial M^{l}[t]}.
\end{equation}
The \emph{choice} of $x(\cdot)$ is crucial—especially when $V^{l}_{\text{thr}}$ is learnable—because it determines both the active region of nonzero gradients and the scaling of their magnitudes.

Reviewing the literature, SG-based methods fall into two families based on how they define $x(M^{l}[t],V^{l}_{\text{thr}})$: \emph{Absolute-Scale} (AS-SG) and \emph{Relative-Scale} (RS-SG). The distinction becomes significant once the threshold is trained.

\section{Two Families of Surrogate Gradients}
\label{sec:trsg_two_families}

To make the scaling effects explicit, we analyze a representative surrogate with a rectangular derivative
\[
f'(x) \;=\; \frac{1}{\gamma}\,\mathds{1}\!\Bigl(|x|<\tfrac{\gamma}{2}\Bigr),
\]
where $\gamma>0$ controls the width of the region where gradients flow. The phenomena described below are \emph{not} specific to this choice; any surrogate exhibits the same issues. \Cref{fig:second_main} visualizes the membrane-potential distribution under two different threshold scales: small ($V_{\text{thr}}^{l}\!\ll\!1$) and large ($V_{\text{thr}}^{l}\!\gg\!1$).

\subsection{AS-SG: Absolute-Scale Surrogate Gradient}
\label{subsec:trsg_assg}

Many works~\cite{zheng:2021, deng:2022, fang:2021, wang:2022} take
\begin{equation}
\label{eq:abs_diff_x}
x \;=\; M^{l}[t] - V^{l}_{\text{thr}},
\end{equation}
which gives
\begin{equation}
\label{eq:abs_diff_grad}
f'\!\bigl(M^{l}[t]-V^{l}_{\text{thr}}\bigr)
\;=\;
\frac{1}{\gamma}\,\mathds{1}\!\Bigl(\bigl|M^{l}[t]-V^{l}_{\text{thr}}\bigr|<\tfrac{\gamma}{2}\Bigr),
\end{equation}
and thus
\begin{equation}
\label{eq:abs_diff_sgrad}
\frac{\partial S^{l}[t]}{\partial M^{l}[t]}
\;\approx\;
f'\!\bigl(M^{l}[t]-V^{l}_{\text{thr}}\bigr).
\end{equation}
Here the active window $\{M^{l}[t]\in(V^{l}_{\text{thr}}-\tfrac{\gamma}{2},\,V^{l}_{\text{thr}}+\tfrac{\gamma}{2})\}$ has a \emph{fixed} width $\gamma$, with magnitude independent of $V^{l}_{\text{thr}}$. When $V^{l}_{\text{thr}}$ is small, the window is excessively wide relative to the membrane-potential spread, so gradients keep flowing far from the true firing boundary—a regime we term \emph{gradient flooding}. Conversely, when $V^{l}_{\text{thr}}$ is large, the fixed-width window covers only a tiny fraction of the distribution, and most gradients vanish—\emph{gradient starvation}.

\subsection{RS-SG: Relative-Scale Surrogate Gradient}
\label{subsec:trsg_rssg}

Other works normalize by the threshold, either by setting $\gamma = V^{l}_{\text{thr}}$~\cite{meng:2023} or by defining
\(
x \;=\; \tfrac{M^{l}[t]}{V^{l}_{\text{thr}}} - 1
\)
as in~\cite{rathi:2023}. In this case,
\begin{equation}
\label{eq:rel_diff_grad}
f'\!\Bigl(\tfrac{M^{l}[t]}{V^{l}_{\text{thr}}}-1\Bigr)
\;=\;
\frac{1}{\gamma}\,\mathds{1}\!\Bigl(\Bigl|\tfrac{M^{l}[t]}{V^{l}_{\text{thr}}}-1\Bigr|<\tfrac{\gamma}{2}\Bigr),
\end{equation}
so the active window $\{M^{l}[t]\in((1-\tfrac{\gamma}{2})V^{l}_{\text{thr}},\,(1+\tfrac{\gamma}{2})V^{l}_{\text{thr}})\}$ scales with $V^{l}_{\text{thr}}$. However, the chain rule introduces a sensitivity of the gradient magnitude to the threshold:
\begin{equation}
\label{eq:rel_diff_sgrad}
\frac{\partial S^{l}[t]}{\partial M^{l}[t]}
\;\approx\;
\frac{1}{V^{l}_{\text{thr}}}\,
f'\!\Bigl(\tfrac{M^{l}[t]}{V^{l}_{\text{thr}}}-1\Bigr).
\end{equation}
Thus, small thresholds amplify gradients excessively (\emph{gradient explosion}), whereas large thresholds suppress them severely (\emph{gradient vanishing}).

\section{Threshold-Robust Surrogate Gradient}
\label{sec:trsg_method}

We propose \textbf{TrSG}, a simple modification that preserves the desirable \emph{relative} window while canceling the problematic $1/V^{l}_{\text{thr}}$ factor in the gradient. During training we (i) keep the relative-scale argument,
\[
x \;=\; \frac{M^{l}[t]}{V^{l}_{\text{thr}}} - 1,
\]
and (ii) multiply the spike by the threshold on the forward path:
\begin{equation}
\label{eq:ta_srgd_core}
O^{l}[t] \;=\; V^{l}_{\text{thr}}\cdot S^{l}[t].
\end{equation}
Then,
\begin{equation}
\label{eq:ta_srgd_chainrule}
\frac{\partial O^{l}[t]}{\partial M^{l}[t]}
\;=\;
V^{l}_{\text{thr}}\;\frac{\partial S^{l}[t]}{\partial M^{l}[t]}
\;\approx\;
V^{l}_{\text{thr}}\cdot
\Bigl(\frac{1}{V^{l}_{\text{thr}}} f'(x)\Bigr)
\;=\;
f'(x),
\end{equation}
so the gradient magnitude becomes \emph{threshold-invariant} while the active window still adapts with $V^{l}_{\text{thr}}$.

\subsection{Inference Remains Binary}
\label{subsec:trsg_inference}

The threshold scaling in \cref{eq:ta_srgd_core} is training-only. At test time we revert to $S^{l}[t]\in\{0,1\}$ and absorb $V^{l}_{\text{thr}}$ into the next layer through a simple weight rescaling
\(
W^{l+1}_{l}\leftarrow V^{l}_{\text{thr}}\,W^{l+1}_{l}.
\)
Inference therefore uses standard binary spikes and a standard LIF neuron, preserving compatibility with existing neuromorphic hardware.

\paragraph{Summary.}
TrSG unifies a relative-scale argument with threshold multiplication on the forward path, ensuring stable gradients across thresholds. This design eliminates both the flooding/starvation issue of AS-SG and the vanishing/explosion problem of RS-SG.

\section{Detailed Gradient Analysis}
\label{sec:trsg_grad_analysis}

In \cref{sec:trsg_two_families,sec:trsg_method} our analysis of threshold-related instabilities centered on $\partial S^{l}[t]/\partial M^{l}[t]$, i.e., how the spike output changes with respect to the membrane potential. However, similar issues arise when considering partial derivatives with respect to the learnable internal parameters $V^{l}_{\text{thr}}$ and $\tau^{l}$. We now derive these gradients explicitly and show that AS-SG and RS-SG fail along these pathways as well, while TrSG remains stable.

\subsection{Gradient with Respect to \texorpdfstring{$V^{l}_{\text{thr}}$}{V\_thr}}
\label{subsec:trsg_grad_vthr}

Recall from \cref{eq:lif_discrete_firing} that $S^{l}[t]=\mathcal{H}\!\bigl(M^{l}[t]-V^{l}_{\text{thr}}\bigr)$. When approximating $\mathcal{H}$ by a differentiable surrogate $f_{sr}$, existing methods define the argument $x$ in either absolute-scale or relative-scale form.

\paragraph{AS-SG.}
With $x = M^{l}[t]-V^{l}_{\text{thr}}$,
\begin{equation}
\frac{\partial S^{l}[t]}{\partial V^{l}_{\text{thr}}}
\;=\;
f_{sr}'(x)\,\Bigl(\tfrac{\partial M^{l}[t]}{\partial V^{l}_{\text{thr}}} - 1\Bigr).
\end{equation}
This update does not adapt its magnitude to $V^{l}_{\text{thr}}$. Consequently, if $|V^{l}_{\text{thr}}|$ becomes very large or very small, the gradient can vanish or dominate excessively.

\paragraph{RS-SG.}
With $x = \tfrac{M^{l}[t]}{V^{l}_{\text{thr}}}-1$, the chain rule yields
\begin{equation}
\frac{\partial S^{l}[t]}{\partial V^{l}_{\text{thr}}}
\;\approx\;
f_{sr}'\!\Bigl(\tfrac{M^{l}[t]}{V^{l}_{\text{thr}}}-1\Bigr)\,\frac{\partial}{\partial V^{l}_{\text{thr}}}\!\Bigl(\tfrac{M^{l}[t]}{V^{l}_{\text{thr}}}-1\Bigr),
\end{equation}
which introduces factors proportional to $\tfrac{1}{V^{l}_{\text{thr}}}$. Tiny $V^{l}_{\text{thr}}$ explodes the gradient; huge $V^{l}_{\text{thr}}$ vanishes it.

\paragraph{TrSG.}
TrSG circumvents both extremes by adopting the relative-scale argument and multiplying $V^{l}_{\text{thr}}$ on the forward path:
\begin{equation}
O^{l}[t] \;=\; V^{l}_{\text{thr}}\,S^{l}[t].
\end{equation}
The resulting gradient of the loss $\mathcal{L}$ with respect to the threshold combines the direct effect of the multiplication and the indirect effect via the surrogate:
\begin{equation}
\frac{\partial \mathcal{L}}{\partial V^{l}_{\text{thr}}}
\;=\;
\frac{\partial \mathcal{L}}{\partial O^{l}[t]}
\Bigl(S^{l}[t] \;+\; V^{l}_{\text{thr}}\,\tfrac{\partial S^{l}[t]}{\partial V^{l}_{\text{thr}}}\Bigr).
\end{equation}
The threshold multiplication cancels the problematic $\tfrac{1}{V^{l}_{\text{thr}}}$ factor in the chain rule, keeping the gradient with respect to $V^{l}_{\text{thr}}$ stable across the full range of threshold scales.

\subsection{Gradient with Respect to \texorpdfstring{$\tau^{l}$}{tau}}
\label{subsec:trsg_grad_tau}

The leakage constant $\tau^{l}$ enters $M^{l}[t]$ through the discrete LIF update of \cref{eq:lif_discrete_potential}:
\begin{equation}
M^{l}[t] \;=\; \Bigl(1-\tfrac{1}{\tau^{l}}\Bigr) U^{l}[t-1] + \tfrac{1}{\tau^{l}} I^{l}_{\text{in}}[t].
\end{equation}
Since $\tau^{l}$ appears only in $M^{l}[t]$, it influences $S^{l}[t]$ indirectly through the chain
\begin{equation}
\frac{\partial S^{l}[t]}{\partial \tau^{l}}
\;=\;
\frac{\partial S^{l}[t]}{\partial M^{l}[t]}
\;\times\;
\frac{\partial M^{l}[t]}{\partial \tau^{l}}.
\end{equation}
If $\partial S^{l}[t]/\partial M^{l}[t]$ is poorly scaled (as in standard AS-SG or RS-SG with extreme $V^{l}_{\text{thr}}$), the corresponding gradient with respect to $\tau^{l}$ can also explode or vanish. By stabilizing the membrane-potential pathway, TrSG simultaneously stabilizes the $\tau^{l}$ pathway, ensuring well-behaved updates of the leakage constant.

\paragraph{Summary.}
TrSG's threshold-multiplication mechanism, combined with its relative-scale argument, ensures that gradients remain appropriately scaled not only along $M^{l}[t]$ but also along the $V^{l}_{\text{thr}}$ and $\tau^{l}$ pathways. This provides a uniformly stable foundation for end-to-end training of LIF SNNs with multiple learnable internal parameters.

\section{Compatibility Across Surrogate Shapes}
\label{sec:trsg_compatibility}

The TrSG construction is independent of the specific surrogate-derivative shape $f'(\cdot)$. We highlight four common choices and confirm that the threshold-invariance property of \cref{eq:ta_srgd_chainrule} carries through.

\paragraph{Rectangular.}
$f'(x) = \tfrac{1}{\gamma}\,\mathds{1}(|x|<\tfrac{\gamma}{2})$.

\paragraph{Triangular.}
$f'(x) = \tfrac{1}{\gamma^{2}}\,\max(0,\gamma-|x|)$.

\paragraph{Arctan.}
$f'(x) = \tfrac{1}{\pi}\,\tfrac{1}{1+(\pi x)^{2}}$.

\paragraph{Sigmoid.}
$f'(x) = \alpha\,\sigma(\alpha x)\,(1-\sigma(\alpha x))$.

For all four, the chain-rule cancellation in \cref{eq:ta_srgd_chainrule} eliminates the $1/V_{\text{thr}}$ factor; \cref{ch:experiments} verifies that empirical training behavior aligns with this prediction.

\section{Summary}
\label{sec:trsg_summary}

This chapter classified existing SG methods into AS-SG and RS-SG, exposed the threshold-induced gradient pathologies of each, and introduced TrSG, which restores threshold invariance with a single forward-path modification. The technique is shape-agnostic, training-only, and inference-compatible with standard LIF neurons.
